\documentclass[11pt,a4paper]{article}
\usepackage[utf8]{inputenc}
\usepackage[T1]{fontenc}
% Tipografia Corporativa Oficial (Helvetica / Helvetica Neue + Courier Mono para datos)
\usepackage[scaled=0.94]{helvet}
\usepackage{courier}
\renewcommand{\familydefault}{\sfdefault}
\usepackage{microtype}
\usepackage{amsmath,amssymb}
\usepackage{booktabs}
\usepackage{tabularx}
\usepackage{colortbl}
\usepackage{enumitem}

% Antigravity Corporate Brand Identity (Interchangeable: bnw, boletos, nodo, vini, wit)
\usepackage[brand=bnw]{../core/agy-corporate-brand}

\geometry{
  top=2.0cm,
  bottom=1.8cm,
  left=1.8cm,
  right=1.8cm,
  headheight=38pt,
  headsep=12pt,
  footskip=16pt
}

% Definición del entorno corptech requerido para documentación técnica SOP
\newtcolorbox{corptech}[2][]{%
  enhanced,
  colback=corpLightBg,
  colframe=corpPrimary,
  coltitle=white,
  fonttitle=\bfseries\sffamily\small,
  arc=2mm,
  boxrule=0.8pt,
  leftrule=4pt,
  title={\textbf{#2}},
  colbacktitle=corpPrimary,
  top=3mm, bottom=3mm, left=4mm, right=4mm,
  drop shadow={opacity=0.05, shadow xshift=1mm, shadow yshift=-1mm},
  #1
}

\begin{document}

% ==================== PÁGINA 1: PORTADA Y DIAGRAMA DE CARRILES (SWIMLANES) ====================
\setupcorporateheader{PROCEDIMIENTO OPERATIVO ESTANDAR (SOP)}{REF: SOP-OPS-2026-042}

\drawchromaticbar[3.5pt]
\vspace{2.5mm}

\noindent
\begin{minipage}[t]{0.65\textwidth}
  {\fontsize{16}{20}\selectfont\bfseries\color{corpPrimary}MANUAL DE PROCESOS OPERATIVOS (SOP)}\\[2pt]
  {\small\bfseries\color{corpTextAccent}GALERIA DE DIAGRAMAS DE FLUJO, ARQUITECTURA \& ESTADOS TRANSACCIONALES}\\[3pt]
  {\footnotesize\color{corpTextMuted}Estandar Corporativo de Diseno de Procesos, Protocolos 2PC y Gobernanza de Sistemas}
\end{minipage}\hfill
\begin{minipage}[t]{0.32\textwidth}
  \raggedleft
  \begin{tcolorbox}[
    enhanced,
    colback=corpLightBg,
    colframe=corpBorder,
    arc=2mm,
    boxrule=0.7pt,
    top=2mm, bottom=2mm, left=3mm, right=3mm
  ]
    \raggedleft
    {\scriptsize\bfseries\color{corpTextMuted}CODIGO DE PROCEDIMIENTO}\\[1pt]
    {\normalsize\bfseries\color{corpPrimary}SOP-OPS-2026-042}\\[2pt]
    {\scriptsize\color{corpTextDark}Version: \textbf{v4.2 Vigente}}\\[1pt]
    {\scriptsize\color{corpSuccess}\textbf{Auditoria: Aprobada ISO 9001}}
  \end{tcolorbox}
\end{minipage}

\vspace{2mm}

\begin{corptech}{Finalidad y Alcance del Procedimiento}{tech-sop-scope}
  Este documento formaliza el flujo integral de adquisicion, validacion fiscal, integracion tecnica y auditoria de infraestructura distribuida. Define las responsabilidades matriciales de Operaciones, Tesoreria, TI y Auditoria, garantizando trazabilidad criptografica en cada etapa bajo acuerdos de nivel de servicio (SLA) corporativos.
\end{corptech}

\vspace{1.5mm}
\noindent{\bfseries\color{corpPrimary}\large 1. Flujo de Proceso Interdepartamental con Carriles (Swimlanes)}
\vspace{1mm}

\noindent\begin{center}
\begin{tikzpicture}[
  scale=0.92, transform shape,
  node distance=1.1cm and 1.2cm,
  swimlane/.style={draw=corpBorder, fill=corpLightBg, rounded corners=2pt, line width=0.6pt},
  laneheader/.style={fill=corpPrimary, text=white, font=\sffamily\bfseries\scriptsize, align=center, inner sep=3pt, minimum width=3.8cm, minimum height=0.6cm}
]
  % Carriles (Swimlanes Columns)
  \draw[swimlane] (0,0) rectangle (4.0, -8.2);
  \draw[swimlane] (4.2,0) rectangle (8.2, -8.2);
  \draw[swimlane] (8.4,0) rectangle (12.4, -8.2);
  \draw[swimlane] (12.6,0) rectangle (16.6, -8.2);

  % Encabezados de Carril
  \node[laneheader] at (2.0, 0.3) {1. OPERACIONES};
  \node[laneheader] at (6.2, 0.3) {2. TESORERIA};
  \node[laneheader] at (10.4, 0.3) {3. TECNOLOGIA / TI};
  \node[laneheader] at (14.6, 0.3) {4. AUDITORIA};

  % Nodos de Proceso
  % Carril 1: Operaciones
  \node[corpevent, minimum width=2.8cm] (start) at (2.0, -0.8) {Inicio de Solicitud};
  \node[corpnode, text width=2.8cm] (op1) at (2.0, -2.1) {Levantamiento de Requerimiento \& OC};
  \node[corpnode, text width=2.8cm] (op2) at (2.0, -6.6) {Recepcion Conforme en Bodega Edge};

  % Carril 2: Tesorería
  \node[corpdecision, text width=2.2cm] (dec_presupuesto) at (6.2, -3.4) {Presupuesto\\Aprobado?};
  \node[corpprimarynode, text width=2.6cm] (tes_reject) at (6.2, -1.8) {Rechazo / Reenvio a Modificacion};
  \node[corpnode, text width=2.6cm] (tes_reserva) at (6.2, -5.0) {Reserva de Fondos \& Emision Fiscal};

  % Carril 3: Tecnología / TI
  \node[corpdecision, text width=2.2cm] (dec_homolog) at (10.4, -5.0) {Homologacion\\Tecnica OK?};
  \node[corpprimarynode, text width=2.6cm] (ti_ticket) at (10.4, -3.4) {Apertura Ticket Soporte N3};
  \node[corpaccentnode, text width=2.6cm] (ti_despliegue) at (10.4, -6.6) {Aprovisionamiento HSM \& Nodo 2PC};

  % Carril 4: Auditoría
  \node[corpnode, text width=2.6cm] (aud_control) at (14.6, -6.6) {Verificacion DTE SIFEN \& Trazabilidad};
  \node[corpevent, fill=corpLightBg, minimum width=2.8cm] (finish) at (14.6, -7.7) {Cierre \& Sellado};

  % Flechas y Enlaces
  \draw[corparrow] (start) -- (op1);
  \draw[corparrow] (op1.east) -- ++(0.6,0) |- (dec_presupuesto.west);
  \draw[corparrow] (dec_presupuesto) -- node[right, font=\scriptsize\bfseries\color{corpDanger}] {No} (tes_reject);
  \draw[corpdottedarrow] (tes_reject.west) -| (op1.north);
  \draw[corparrow] (dec_presupuesto) -- node[right, font=\scriptsize\bfseries\color{corpSuccess}] {Si} (tes_reserva);
  \draw[corparrow] (tes_reserva) -- (dec_homolog);
  \draw[corparrow] (dec_homolog) -- node[right, font=\scriptsize\bfseries\color{corpWarning}] {No} (ti_ticket);
  \draw[corpdottedarrow] (ti_ticket.west) -- (tes_reserva.north);
  \draw[corparrow] (dec_homolog) -- node[right, font=\scriptsize\bfseries\color{corpSuccess}] {Si} (ti_despliegue);
  \draw[corparrow] (ti_despliegue.west) -- (op2.east);
  \draw[corparrow] (ti_despliegue.east) -- (aud_control.west);
  \draw[corparrow] (aud_control) -- (finish);

\end{tikzpicture}
\end{center}

\clearpage

% ==================== PÁGINA 2: MÁQUINA DE ESTADOS 2PC ====================
\setupcorporateheader{PROCEDIMIENTO OPERATIVO ESTANDAR (SOP)}{REF: SOP-OPS-2026-042}

\drawchromaticbar[3.5pt]
\vspace{2.5mm}

\noindent{\bfseries\color{corpPrimary}\large 2. Maquina de Estados Finitos de Transaccion Criptografica (Two-Phase Commit -- 2PC)}
\vspace{1.5mm}

\begin{corpnote}{Definicion del Protocolo 2PC Distribuido}{note-2pc-protocol}
  Para transacciones de cobro en terminales Boletos, telemetria VINI e integracion tributaria SIFEN / DTE, el coordinador ejecuta un protocolo de compromiso en dos fases con garantia ACID y recuperacion ante caidas sin perdida de sincronia.
\end{corpnote}

\vspace{3mm}

\noindent\begin{center}
\begin{tikzpicture}[
  scale=0.95, transform shape,
  node distance=2.2cm and 2.5cm,
  state/.style={rectangle, rounded corners=6pt, draw=corpPrimary, line width=1.2pt, fill=white, text=corpTextDark, font=\sffamily\footnotesize\bfseries, align=center, inner sep=8pt, drop shadow={opacity=0.08, shadow xshift=1pt, shadow yshift=-1pt}},
  coordstate/.style={rectangle, rounded corners=6pt, draw=corpTextAccent, line width=1.2pt, fill=corpLightBg, text=corpPrimary, font=\sffamily\footnotesize\bfseries, align=center, inner sep=8pt},
  succstate/.style={rectangle, rounded corners=6pt, draw=corpSuccess, line width=1.5pt, fill=corpLightBg, text=corpSuccess, font=\sffamily\footnotesize\bfseries, align=center, inner sep=8pt},
  failstate/.style={rectangle, rounded corners=6pt, draw=corpDanger, line width=1.5pt, fill=white, text=corpDanger, font=\sffamily\footnotesize\bfseries, align=center, inner sep=8pt},
  trtext/.style={font=\scriptsize\sffamily, text=corpTextDark, align=center, fill=white, inner sep=2pt, rounded corners=1pt}
]

  % Nodos de Estado
  \node[coordstate, minimum width=3.2cm] (init) at (0, 0) {ESTADO INIT\\{\scriptsize [Transaccion Creada]}};
  \node[state, minimum width=3.2cm] (prepare) at (5.0, 0) {FASE 1: PREPARE\\{\scriptsize [Bloqueo de Recursos]}};
  \node[state, minimum width=3.2cm] (prepared) at (10.5, 0) {ESTADO PREPARED\\{\scriptsize [Votos Recibidos 100\%]}};
  
  \node[succstate, minimum width=3.4cm] (commit) at (10.5, -4.5) {FASE 2: COMMIT GLOBAL\\{\scriptsize [Persistencia Confirmada]}};
  \node[failstate, minimum width=3.4cm] (abort) at (5.0, -4.5) {FASE 2: ROLLBACK / ABORT\\{\scriptsize [Liberacion Inmediata]}};

  % Transiciones y Flechas
  \draw[corparrow] (init) -- node[above, trtext] {1. tx\_begin()\\Coordinador} (prepare);
  \draw[corparrow] (prepare) -- node[above, trtext] {2. VOTE\_COMMIT\\Nodos OK} (prepared);
  \draw[corparrow] (prepared) -- node[right, trtext] {3. Global Commit\\Broadcast} (commit);
  
  % Rama de Falla / Timeout / Veto
  \draw[corparrow, draw=corpDanger] (prepare) -- node[right, trtext, text=corpDanger] {Timeout /\\VOTE\_ABORT} (abort);
  \draw[corparrow, draw=corpDanger] (prepared) -- node[above, trtext, text=corpDanger] {Falla en Nodo Participante} (abort);
  
  % Ciclo de reintento / Cierre
  \draw[corpdottedarrow] (abort.west) -| node[near start, below, trtext] {Reintento Exponencial (Max 3)} (init.south);
  \draw[corpaccentarrow] (commit.east) -- ++(1.5,0) |- node[near end, above, trtext] {Auditoria SHA-256 Exitosa} (14, 2) -| (init.north);

\end{tikzpicture}
\end{center}

\vspace{3mm}

% Tabla de Matriz de Estados y Códigos de Error
\noindent{\bfseries\color{corpPrimary}\small Matriz de Transiciones y Respuestas del Servidor Coordinador}
\vspace{1mm}

\noindent
\begin{tabularx}{\linewidth}{@{} l l l X @{}}
  \toprule
  \rowcolor{corpPrimary}
  {\bfseries\color{white}\scriptsize Estado Origen} &
  {\bfseries\color{white}\scriptsize Mensaje / Estimulo} &
  {\bfseries\color{white}\scriptsize Estado Destino} &
  {\bfseries\color{white}\scriptsize Accion Operativa y Garantia de Integridad} \\
  \midrule
  \rowcolor{corpLightBg}
  INIT & PREPARE\_REQ & PREPARE & Escritura en log transaccional Write-Ahead (WAL). \\
  PREPARE & ALL\_VOTES\_COMMIT & PREPARED & Verificacion de quorums en Nodos Edge NODO-X8. \\
  \rowcolor{corpLightBg}
  PREPARED & GLOBAL\_COMMIT & COMMIT & Cierre fiscal DTE/SIFEN y emision de comprobante. \\
  PREPARE / PREPARED & TIMEOUT\_OR\_FAIL & ROLLBACK & Deshacer bloqueos y notificar alarma de red N3. \\
  \bottomrule
\end{tabularx}

\clearpage

% ==================== PÁGINA 3: ARQUITECTURA DE SISTEMAS Y CRONOGRAMA ====================
\setupcorporateheader{PROCEDIMIENTO OPERATIVO ESTANDAR (SOP)}{REF: SOP-OPS-2026-042}

\drawchromaticbar[3.5pt]
\vspace{2.5mm}

\noindent{\bfseries\color{corpPrimary}\large 3. Diagrama de Arquitectura de Sistemas Multinivel (Edge to Datacenter)}
\vspace{1mm}

\noindent\begin{center}
\begin{tikzpicture}[
  scale=0.92, transform shape,
  layerbox/.style={draw=corpBorder, fill=corpLightBg, rounded corners=4pt, line width=0.8pt, inner sep=6pt},
  layertitle/.style={font=\sffamily\bfseries\scriptsize, text=corpPrimary, anchor=north west},
  archnode/.style={rectangle, rounded corners=3pt, draw=corpPrimary, fill=white, font=\sffamily\scriptsize, inner sep=5pt, align=center, line width=0.8pt}
]

  % Capa 1: Edge & Dispositivos de Campo
  \node[layerbox, minimum width=15.8cm, minimum height=1.6cm] (l1) at (7.9, 0) {};
  \node[layertitle] at (0.2, 0.7) {CAPA 1: DISPOSITIVOS EDGE \& TERMINALES DE CAMPO (VINI / BOLETOS)};
  \node[archnode] (d1) at (2.2, 0) {Terminal POS Pasajes\\Modem Dual LTE};
  \node[archnode] (d2) at (6.0, 0) {Modulo CAN VINI V4\\Telemetria Flotas};
  \node[archnode] (d3) at (9.8, 0) {Sensoristica IoT\\Monitoreo Activos};
  \node[archnode] (d4) at (13.6, 0) {Validador QR Torniquete\\Sub-segundo};

  % Capa 2: Gateway y Seguridad Perimetral
  \node[layerbox, minimum width=15.8cm, minimum height=1.6cm] (l2) at (7.9, -2.3) {};
  \node[layertitle] at (0.2, -1.6) {CAPA 2: GATEWAY DE BORDE, SEGURIDAD PERIMETRAL \& VPN (NODO-EDGE)};
  \node[corpprimarynode, font=\sffamily\scriptsize] (g1) at (3.5, -2.3) {Gateway Nodo-X8\\Balanceo \& Cifrado WireGuard};
  \node[corpprimarynode, font=\sffamily\scriptsize] (g2) at (8.0, -2.3) {Firewall Layer 7 UTM\\DDoS Protection \& Rate Limit};
  \node[corpprimarynode, font=\sffamily\scriptsize] (g3) at (12.5, -2.3) {Buffer de Contingencia Offline\\Store \& Forward Seguro};

  % Capa 3: Servicios y Lógica Transaccional
  \node[layerbox, minimum width=15.8cm, minimum height=1.6cm] (l3) at (7.9, -4.6) {};
  \node[layertitle] at (0.2, -3.9) {CAPA 3: SERVICIOS CORE, MOTOR 2PC \& CONECTIVIDAD FISCAL};
  \node[corpaccentnode, font=\sffamily\scriptsize] (s1) at (3.2, -4.6) {Coordinador 2PC\\Cluster Kubernetes};
  \node[corpaccentnode, font=\sffamily\scriptsize] (s2) at (7.9, -4.6) {Motor Fiscal DTE / SIFEN\\e-Kuatia WebService REST};
  \node[corpaccentnode, font=\sffamily\scriptsize] (s3) at (12.6, -4.6) {Plataforma Telemetria\\Ingestion Apache Kafka};

  % Capa 4: Datacenter y Persistencia
  \node[layerbox, minimum width=15.8cm, minimum height=1.6cm] (l4) at (7.9, -6.9) {};
  \node[layertitle] at (0.2, -6.2) {CAPA 4: ALMACENAMIENTO CRIPTOGRAFICO \& DATACENTER ENTERPRISE};
  \node[archnode] (db1) at (3.5, -6.9) {PostgreSQL HA Distributed\\Replicacion Sincrona};
  \node[archnode] (db2) at (8.0, -6.9) {Modulo HSM Hardware\\Sellado SHA-256 \& Claves PKI};
  \node[archnode] (db3) at (12.5, -6.9) {Cold Storage / Auditoria\\WORM Compliance 7 Anos};

  % Flechas entre capas
  \draw[corparrow] (d1) -- (g1);
  \draw[corparrow] (d2) -- (g1);
  \draw[corparrow] (d3) -- (g2);
  \draw[corparrow] (d4) -- (g3);

  \draw[corparrow] (g1) -- (s1);
  \draw[corparrow] (g2) -- (s2);
  \draw[corparrow] (g3) -- (s3);

  \draw[corparrow] (s1) -- (db1);
  \draw[corparrow] (s2) -- (db2);
  \draw[corparrow] (s3) -- (db3);

\end{tikzpicture}
\end{center}

\vspace{1.5mm}
\noindent{\bfseries\color{corpPrimary}\large 4. Cronograma de Implementacion y Roadmap de Fases}
\vspace{1mm}

\noindent\begin{center}
\begin{tikzpicture}[
  scale=0.92, transform shape,
  phase/.style={rectangle, rounded corners=2pt, font=\sffamily\scriptsize\bfseries, align=center, inner sep=4pt},
  mile/.style={diamond, aspect=1.2, fill=corpPrimary, text=white, font=\sffamily\tiny\bfseries, inner sep=2pt}
]
  % Eje de Tiempo (Trimestres)
  \draw[-Latex, line width=1.2pt, color=corpPrimary] (0, 0) -- (16.2, 0);
  \node[font=\sffamily\scriptsize\bfseries\color{corpPrimary}, anchor=south] at (2.0, 0.1) {Q1-2027: Diseno};
  \node[font=\sffamily\scriptsize\bfseries\color{corpPrimary}, anchor=south] at (6.0, 0.1) {Q2-2027: Homologacion};
  \node[font=\sffamily\scriptsize\bfseries\color{corpPrimary}, anchor=south] at (10.0, 0.1) {Q3-2027: Despliegue};
  \node[font=\sffamily\scriptsize\bfseries\color{corpPrimary}, anchor=south] at (14.0, 0.1) {Q4-2027: Certificacion};

  % Bloques de Fase
  \node[phase, fill=corpPrimary, text=white, minimum width=3.4cm, minimum height=0.6cm] at (2.0, -0.8) {Fase I: Arquitectura Edge};
  \node[phase, fill=corpTextAccent, text=white, minimum width=3.4cm, minimum height=0.6cm] at (6.0, -1.6) {Fase II: Piloto SIFEN / DTE};
  \node[phase, fill=corpSecondary, text=white, minimum width=3.4cm, minimum height=0.6cm] at (10.0, -0.8) {Fase III: Rollout 120 Nodos};
  \node[phase, fill=corpSuccessGreen, text=white, minimum width=3.4cm, minimum height=0.6cm] at (14.0, -1.6) {Fase IV: Auditoria ISO 27001};

  % Hitos (Milestones)
  \node[mile] (m1) at (3.8, 0) {H1};
  \node[mile] (m2) at (7.8, 0) {H2};
  \node[mile] (m3) at (11.8, 0) {H3};
  \node[mile] (m4) at (15.8, 0) {H4};

  \node[font=\tiny\sffamily\color{corpTextMuted}, below=2pt of m1] {Prototipo OK};
  \node[font=\tiny\sffamily\color{corpTextMuted}, below=2pt of m2] {Cert. Fiscal};
  \node[font=\tiny\sffamily\color{corpTextMuted}, below=2pt of m3] {Cluster 2PC};
  \node[font=\tiny\sffamily\color{corpTextMuted}, below=2pt of m4] {Go-Live 100\%};
\end{tikzpicture}
\end{center}

\vspace{1.5mm}

% Firmas de Aprobación Técnica y Operativa
\dualcontrolsignature%
  {Ing. Rodrigo Morales P.}{Gerente Corporativo de Infraestructura \& Redes}{Aprobacion Tecnica Arquitectura}%
  {Dra. Claudia Miranda Valenzuela}{Vicepresidente de Auditoria \& Cumplimiento}{Conformidad Procedimiento SOP}

\vspace{2mm}

% Sello Criptográfico
\cryptoseal%
  {9a74b3f110c83e45dc56187b89e248b1d98a0928a6fbb34e9087541289cfda01}%
  {2026-11-03T12:50:00.000Z}%
  {SOP-TECH-PIPELINE-VALIDATED-NODE-SHA256}

\end{document}
